\IfFileExists{IEEEtran.cls}{
  \documentclass[conference]{IEEEtran}
}{
  \documentclass[twocolumn]{article}
  \newcommand{\IEEEauthorblockN}[1]{\textbf{##1}}
  \newcommand{\IEEEauthorblockA}[1]{##1}
}

\usepackage{amsmath,amssymb,amsthm}
\usepackage{mathtools}
\usepackage{booktabs}
\usepackage{graphicx}
\usepackage{hyperref}
\usepackage{xurl}
\usepackage{microtype}
\usepackage{enumitem}
\usepackage{placeins}
\usepackage{pgfplots}
\pgfplotsset{compat=1.18}
\providecommand{\appendices}{\appendix}

\newcommand{\PaperTitle}{Certificate-Guided Decomposition for Explainable and Repairable University Timetabling}
\newcommand{\PaperAuthor}{Mahmoud Elfeel}
\newcommand{\PaperAffiliation}{Independent Timetabling Study}
\newcommand{\PaperEmail}{mahmoudelfeelig@gmail.com}

\hypersetup{
  colorlinks=true,
  linkcolor=black,
  citecolor=black,
  urlcolor=blue,
  pdftitle={\PaperTitle},
  pdfauthor={\PaperAuthor},
  pdfsubject={University course timetabling, room decomposition, explanations, and repair},
  pdfkeywords={course timetabling, CP-SAT, Benders decomposition, Hall certificate, fairness, repair}
}
\title{\PaperTitle}
\author{
\IEEEauthorblockN{\PaperAuthor}
\IEEEauthorblockA{
\textit{\PaperAffiliation} \\
\PaperEmail
}
}

\begin{document}
\maketitle

\begin{abstract}
Time/room decomposition is an established approach to university course timetabling, but operational systems also need to explain why a fixed-time room assignment fails and identify a small, safe repair scope.
We present a research prototype that connects these concerns through certificate-guided decomposition.
A CP-SAT master assigns activity times; an exact room subproblem either returns a room assignment or emits a Hall-deficiency certificate for a capacity- and availability-sensitive room neighborhood.
The certificate produces a standard option-expanded Hall cut with right-hand side $|\Gamma(U)|$, a human-readable counting explanation, and a resource-linked large-neighborhood repair scope.
A separate replay checker reconstructs the effective domains and cut before insertion, while content-addressed SHA-256 identifiers link certificate-derived neighborhood traces as integrity metadata rather than signatures.
Non-Hall failures caused by travel, repeat-room, or lock interactions yield exact schedule nogoods.
The implementation additionally provides adaptive parallel week decomposition with explicit cross-partition safeguards, scenario-robust room demand, portable institutional distribution constraints, executable policy presets, an exact lexicographic fairness profile, deterministic experiment seeds, official ITC-2007 soft scoring, disruption trials, and a checksum-verifiable artifact workflow.
We distinguish implementation verification from empirical evidence: earlier generated tables are quarantined because their seed labels did not change instances and their strict-room timings included retry/fallback work.
A replacement one-worker study preserves 150 primary runs over 50 paired generated instances and passes a preregistered minimum of 30 unique feasible instances per condition under a separate repository-local post-solve validator; exact decomposition reduces paired time-to-feasible relative to the integrated room model, while greedy rooming remains faster.
On the 21 official ITC-2007 instances, two fresh 42-row passes are fully official-validator feasible: Planora/CPSolver/tie counts are 9/10/2 and 14/7/0, and Planora has the lower mean objective in both passes, but seven instance winners change between repeats.
A corrected, counterbalanced, single-seed fixed-time room-dive ablation is official-validator feasible and deadline-compliant in all 42 runs and lowers the ON-condition aggregate by 2.869\%; however, only one accepted dive supplies a direct within-run gain, and the feature remains disabled by default.
The generated result is limited to feasibility latency, while repeat sensitivity still rules out a general superiority claim. Authenticated ITC-2019 validation agrees with every complete locally valid output produced by a bounded four-solver matrix, but only three of 120 competition conditions emit such an output; effective breadth, certificate-guided repair benefit, multi-seed external quality, and institutional sign-off remain open gates.
\end{abstract}

\section{Introduction}
University course timetabling must place activities in time and rooms while preserving student, staff, equipment, capacity, and institutional constraints.
The optimization problem is mature, but the operational failure mode remains difficult: a time-feasible schedule can still be impossible to room, and a bare ``infeasible'' status does not tell an administrator which activities must move.

Time/room decomposition is not new.
Bagger et al.\ explicitly decompose curriculum-based course timetabling into time scheduling and room allocation~\cite{bagger2018benders}; hybrid and metaheuristic approaches are also well established~\cite{muller2009hybrid,lewis2008metaheuristics,bettinelli2015overview}.
Fair timetabling based on max--min orderings and Jain's index has likewise been studied~\cite{muhlenthaler2013fairness,muhlenthaler2013decomposition}.
The research question in this work is therefore not whether decomposition or fairness works in general.
It is whether a room-subproblem certificate can serve three roles in one auditable system: a valid master cut, a concise explanation, and a bounded repair neighborhood.

The implemented contributions are:
\begin{itemize}[leftmargin=*]
  \item a CP-SAT time master coupled to an exact, full-domain room assignment subproblem;
  \item Hall-deficiency extraction for simultaneous room jobs, with capacity, room type, specialization, availability, closure, and lock filtering performed before the certificate is formed;
  \item iterative certificate cuts that remove a deficient occupancy pattern, plus exact nogoods for infeasibility caused by travel, repeat-room, or lock interactions not captured by a slot-wise Hall witness;
  \item reuse of the same certificate for human-readable explanations and a resource-linked repair neighborhood while freezing unaffected activities;
  \item an adaptive large-instance path that validates global invariants, solves decision-independent weeks concurrently, accepts only validated room assignments, activates exact certificate decomposition on a failing week, and refuses hard cross-week coupling rather than silently weakening it;
  \item enrollment-uncertainty policies whose nominal, scenario-quantile, worst-case, or budgeted demand is reused consistently by room domains, validators, metrics, and infeasibility certificates;
  \item a portable distribution-constraint schema with exact evaluation for nineteen common ITC-style relations, hard CP compilation for seventeen of them, and fail-fast handling instead of silent weakening for the two evaluator-only aggregate rules;
  \item executable institution-policy presets, including a portable baseline and a repository-derived GIU target preset that is explicitly marked as requiring local policy validation;
  \item an exact lexicographic fairness profile that minimizes the maximum group burden before total weighted penalty, accompanied by tail, Gini, and Jain diagnostics;
  \item a reproducible evaluation harness with effective instance/solver/local-search seeds, instance and schedule fingerprints, solver bounds and gaps, paired statistics, official ITC-2007 soft scoring, ITC-2019 schema inspection, portable metrics, disruption experiments, and verified release bundles.
\end{itemize}

The systematic review falsifies the broad mathematical novelty claim: the contextual cut is an option-expanded instance of the established partial-transversal Hall system, while Benders-guided and explanation-guided LNS both have direct prior art.
The remaining research hypothesis is a systems and empirical one: whether exact effective-domain reconstruction, conservative cut fallback, typed certificate-support transfer, reusable exact repair, and auditable lineage improve performance and operational trust when implemented together.
A publication submission must still complete the evaluation protocol in Section~VII and an independent literature review.
This paper deliberately separates implemented capability, verified tests, planned experiments, and unsupported claims.


\section{Related Work}
Automated timetabling spans mathematical programming, constraint programming, local search, and hybrids~\cite{schaerf1999survey,lewis2008metaheuristics,bettinelli2015overview}.
Competition formats made hard/soft semantics, validators, and comparable instances central to the field~\cite{mccollum2010itc,bonutti2012benchmarking}.
ITC~2019 broadened the target to real-world classes, admissible time/room domains, distribution constraints, and student sectioning~\cite{muller2025itc2019}.
MaxSAT preprocessing followed by local improvement is an established ITC~2019 competitor design~\cite{lemos2022unicort}.
Recent surveys emphasize that institution-specific rules and reproducible benchmark validation remain central obstacles to transferring one timetabling formulation between universities~\cite{rossi2022survey}.

\paragraph{Decomposition.}
Logic-based Benders methods use a master problem and a subordinate inference problem whose failures generate cuts~\cite{hooker2003logic}.
Lach and L\"ubbecke characterize the partial-transversal structure behind time-first, room-second university timetabling through Hall inequalities~\cite{lach2008partial}; Bagger et al. segment the curriculum-based problem into time scheduling and room allocation and generate Benders feasibility cuts~\cite{bagger2018benders}.
Consequently, neither time/room decomposition nor Hall/Benders room cuts are claimed as novel.
Expanding the left side of the bipartite graph from jobs to admissible job--start options also subsumes the start-dependent contextual formula: for option set $U$, the established partial-transversal inequality is $x(U)\le |\Gamma(U)|$.
An earlier reviewed draft used the witnessed-room cardinality and could be weaker when $\Gamma(U)$ was a strict subset of that witness; the implementation now uses the stronger standard right-hand side $|\Gamma(U)|$.
This correction strengthens the separation but does not change the prior-art verdict: the method is evaluated as a dynamic effective-domain instantiation with conservative fallback, not as a new inequality family.

\paragraph{Hybrid improvement.}
Feasibility construction followed by tabu search, simulated annealing, variable-neighborhood search, or related neighborhoods is established practice~\cite{lewis2008metaheuristics,muller2009hybrid}.
Burke et al. explicitly combine decomposition, reformulation, and diving for university course timetabling, while L\"u and Hao construct periods heuristically and solve the induced room assignment as an exact bipartite matching problem~\cite{burke2010diving,lu2008hybrid}.
Accordingly, fixing event times and reoptimizing rooms is an established special case rather than a new algorithmic primitive.
Instance-feature-driven parameter tuning is also established for curriculum-based timetabling~\cite{bellio2016feature}; any future learned or calibrated controller must be evaluated against that line of work rather than presented as a new idea in isolation.
CP-based exact repair in course-timetabling LNS, failure-derived cuts, adaptive neighborhood size, Benders-cut-guided LNS, Benders-aware LNS subproblems, and explanation-derived neighborhoods also precede this work~\cite{cambazard2012cp,nagata2018random,trick2011benderslns,maher2021benderslns,prudhomme2014explanationlns}.
Adaptive LNS and online bandit selection are established across mixed-integer and scheduling domains~\cite{ropke2006alns,hendel2022alns,phan2024balance,cai2025balans}.
The Planora local search, exact repair, and discounted-UCB controller are therefore not presented as new metaheuristic classes.
The testable systems hypothesis is whether a typed certificate-support record, with independently replayable cut assumptions and explicit neighborhood lineage, can safely serve as a cut, semantic repair seed, explanation, and audit event.
Room-majority guidance and course room-stability penalties likewise do not define a new neighborhood.
The support-aware feedback ablation evaluated here changes only the deterministic secondary ranking signal: it estimates the number of distinct courses fragmented after a period--room collision projection, while the official four-component score and independent validator retain sole acceptance authority.
It is therefore described as a Planora selection hypothesis, not as a new or exact room-stability objective.

\paragraph{Systematic novelty boundary.}
The reproducible search protocol and prior-art matrix are provided in \path{paper/novelty_review.md} and its machine-readable companion.
An adversarial second pass found direct mathematical and algorithmic anticipation for both broad candidate claims.
We therefore make no priority claim for the contextual inequality, Benders/certificate-guided LNS, or assumption-core-guided repair.
The fixed-time room dive and feature-based tuning directions are likewise prior art.
The remaining candidate is an integration and reproducibility contribution, and it still requires independent database review before submission.

\paragraph{Fairness.}
M\"uhlenthaler and Wanka formulate academic timetabling with max--min fairness and Jain's index, and later study a decomposition of the max--min fair room assignment problem~\cite{muhlenthaler2013fairness,muhlenthaler2013decomposition}.
Our fairness profile therefore implements an established normative direction: minimize the worst group burden lexicographically before aggregate penalty.
Gini and Jain values are diagnostics, not proof that the chosen burden definition is institutionally fair.

\paragraph{Explainability and repair.}
An infeasibility explanation is useful only if its semantics match the actual solver domains.
Our Hall witness is generated after applying type, combined cluster capacity, specialization, availability, closure, and room-lock filters.
Travel and cross-week repeat-room coupling can make the full room model infeasible even when every slot-wise bipartite graph has a perfect matching; those cases receive an exact time-assignment nogood rather than a false Hall explanation.
This distinction is central to the proposed evaluation.

\paragraph{Robustness and institutional policy.}
Enrollment and disruption uncertainty motivate robust or stochastic scheduling rather than deterministic capacity assumptions alone~\cite{holm2023robust}.
Published university policies additionally expose operational criteria not represented by a single competition score, including standard meeting patterns, prime-time balancing, historical enrollment, equipment and location fit, room-fill targets, and post-registration change controls~\cite{ucdavispolicy,berkeleypolicy,michiganpolicy}.
Planora represents these criteria as data and metrics; it does not claim that one generic preset is an official policy for any named institution.

\paragraph{Scale decomposition.}
Institutional systems may stage or partition scheduling work by department, lecture size, laboratory category, or other resource boundaries; UniTime documents a Purdue workflow that first handles large lectures, then departments, and then computing laboratories~\cite{unitimeproblem}.
This supports decomposition as an operational pattern, not a novelty claim.
Planora's additional research question is narrower: whether a policy-completeness preflight plus adaptive certificate fallback can retain auditability while exploiting genuinely independent temporal partitions.


\section{Problem Definition and Notation}
We consider a discrete-time timetabling problem on a weekly grid.
The entity/constraint structure is aligned with curriculum-based course timetabling conventions used in ITC-style benchmarks~\cite{bonutti2012benchmarking,bettinelli2015overview}.

\paragraph{Calendar.}
Let $\mathcal{W}$ be the set of teaching weeks, $\mathcal{D}$ the set of teaching days, and $\mathcal{S}=\{0,\dots,S-1\}$ the set of within-day time slots.
We assume a fixed set of teaching days (Monday--Saturday) and exclude Sunday.

\paragraph{Entities.}
We are given:
\begin{itemize}[leftmargin=*]
  \item groups $g \in \mathcal{G}$ (student cohorts), each with a size $\mathrm{size}(g)$,
  \item staff $p \in \mathcal{P}$ (professors and TAs), each with
    day-availability $\mathrm{AvailDays}(p)\subseteq \mathcal{D}$ and optional
    week-availability $\mathrm{AvailWeeks}(p)\subseteq\mathcal{W}$,
  \item rooms $r \in \mathcal{R}$ with type $\mathrm{type}(r)$ and capacity $\mathrm{cap}(r)$,
    where capacity can be configured either categorically (\texttt{SMALL/MEDIUM/BIG})
    or numerically (exact seat counts),
  \item activities $a \in \mathcal{A}$, each with:
    \begin{itemize}[leftmargin=*]
      \item week $\mathrm{week}(a)\in\mathcal{W}$,
      \item kind $\mathrm{kind}(a)\in\{\mathrm{LEC},\mathrm{TUT},\mathrm{LAB}\}$,
      \item duration $\mathrm{dur}(a)\in\{1,2,3\}$ (number of consecutive slots),
      \item participating groups $\mathcal{G}(a)\subseteq \mathcal{G}$,
      \item assigned staff $\mathrm{staff}(a)\in\mathcal{P}$,
      \item optional specialization tag requirement $\mathrm{tag}(a)$ for specialized labs.
    \end{itemize}
\end{itemize}

\paragraph{Decision.}
For each activity $a$, choose a start day and start slot $(d,s)$ in week $\mathrm{week}(a)$; optionally choose a room $r$.

\paragraph{Feasibility.}
Hard constraints include resource non-overlap (groups, staff, rooms), staff availability, room eligibility, and several rule-based constraints such as ``week 1 is lectures-only'' and optional staff load caps.

\paragraph{Quality.}
We also consider soft constraints expressed as a weighted penalty:
free-day shortfalls, gaps, late starts, stability, room consistency,
and an additional weekly concentration term discouraging many repeated
same-course LEC/TUT sessions for one group in one week (useful during manual week-shift editing).


\section{CP-SAT Formulation}
This section presents the scheduling formulation in a standard mixed 0--1 form.
CP-SAT is used as the optimization engine, with interval variables for non-overlap constraints.

\subsection{Time Encoding and Start Variables}
For each week, index time within the week by $t \in \{0,\dots,|\mathcal{D}|S-1\}$ where
$t = \mathrm{dayIndex}(d)\cdot S + s$.

Each activity $a$ has a set of allowed start indices $T_a$ derived from staff
day/week availability, duration feasibility, institutional windows, generic-resource availability, and optional locks.
The implementation uses one integer variable $s_a\in T_a$ rather than eagerly materializing a one-hot vector for every domain value.
When a constraint, objective term, or decomposition cut requires a value literal, $x_{a,t}\Leftrightarrow(s_a=t)$ is created lazily.
Objective-enabled models may create the full one-hot channeling once because many occupancy penalties reuse it:
\begin{align}
\sum_{t\in T_a} x_{a,t} &= 1, &
s_a &= \sum_{t\in T_a}t x_{a,t}. \label{eq:oneStart}
\end{align}
Pure feasibility models therefore scale with one start integer per activity plus only the literals demanded by active rules and cuts.

\paragraph{Staff week availability.}
If $\mathrm{week}(a)\notin \mathrm{AvailWeeks}(\mathrm{staff}(a))$, then
$T_a=\emptyset$ and the instance is infeasible under that assignment policy.
Equivalently, in feasibility checks and manual edits we enforce:
\[
\mathrm{week}(a)\in \mathrm{AvailWeeks}(\mathrm{staff}(a)).
\]

\subsection{Interval Non-Overlap Constraints}
Let $\mathrm{Occ}(a,t,\tau)$ be 1 iff slot $\tau$ is covered when activity $a$ starts at $t$, i.e.,
$\tau \in [t, t+\mathrm{dur}(a)-1]$.

\paragraph{Group conflicts.}
For each group $g$, week $w$, and time slot $\tau$:
\begin{align}
\sum_{\substack{a\in\mathcal{A}:\\ g\in\mathcal{G}(a),\,\mathrm{week}(a)=w}}
\;\sum_{t\in T_a:\,\mathrm{Occ}(a,t,\tau)=1} x_{a,t} \;\le\; 1.
\end{align}

\paragraph{Staff conflicts.}
Analogously, for each staff member $p$, week $w$, and time slot $\tau$:
\begin{align}
\sum_{\substack{a\in\mathcal{A}:\\ \mathrm{staff}(a)=p,\,\mathrm{week}(a)=w}}
\;\sum_{t\in T_a:\,\mathrm{Occ}(a,t,\tau)=1} x_{a,t} \;\le\; 1.
\end{align}

The implementation uses CP-SAT interval variables and \texttt{AddNoOverlap} for these constraints.

\subsection{Rule Constraints}
\paragraph{Week-1 lectures-only.}
If $w_0=\min(\mathcal{W})$, then activities of kind TUT/LAB are forbidden in week $w_0$.
The implementation rejects such instances during model construction.

\paragraph{Load caps (optional).}
If $\mathrm{MaxWeek}(p)$ is specified, enforce:
\begin{align}
\sum_{\substack{a:\,\mathrm{staff}(a)=p,\\\mathrm{week}(a)=w}}
\mathrm{dur}(a) &\le \mathrm{MaxWeek}(p),
\qquad \forall p,w.
\end{align}
Similarly, if $\mathrm{MaxDay}(p)$ is specified, enforce per-day caps.

\paragraph{Block-only professors.}
For staff with a \emph{block-only} flag, lectures for a given course in a given week must form a single contiguous block of length 2--3 on one day.
The implementation encodes this via per-slot occupancy indicators and a ``single block-start'' constraint.

\subsection{Institutional Distribution Constraints}
An instance can carry typed relations over an ordered activity set.
The shared evaluator implements \texttt{SameStart}, \texttt{SameTime}, \texttt{DifferentTime}, day/week equality and inequality, overlap/non-overlap, room equality and inequality, same attendees, precedence, work-day span, minimum gap, maximum days, maximum day load, maximum breaks, and maximum block.
Hard pair relations are compiled as allowed-domain tables; room relations are compiled either into integrated room selectors or the exact room subproblem.
Maximum-day and maximum-day-load rules use day-presence literals.
Hard maximum-break and maximum-block rules currently fail at model construction because only their exact evaluator and soft scorer are implemented; this explicit capability boundary prevents an imported rule from being weakened silently.

\subsection{Enrollment-Uncertainty Capacity}
Let nominal demand for activity or cluster $J$ be $\bar d_J=\sum_{g\in G(J)}\bar d_g$.
Named forecast scenario $\omega$ gives $d_J^\omega=\sum_g d_g^\omega$.
Let $k=\lfloor\Gamma\rfloor$ and define budgeted demand as
\begin{align}
D_J^{\mathrm{bud}}
=\bar d_J+\sum_{i=1}^{k}\hat d_{(i)}
+\left\lceil(\Gamma-k)\hat d_{(k+1)}\right\rceil .
\end{align}
The selectable room requirement is
\begin{align}
D_J = \begin{cases}
\bar d_J & \text{nominal},\\
\max_\omega d_J^\omega & \text{worst case},\\
Q_q(\{d_J^\omega\}) & \text{empirical quantile }q,\\
D_J^{\mathrm{bud}} & \text{budgeted},
\end{cases}
\end{align}
where $\hat d_{(i)}$ are descending group deviations.
The same $D_J$ is used by integrated room domains, exact room jobs, greedy eligibility, post-solve validation, feasibility diagnostics, and portable service metrics.

\subsection{Clustering (Co-Location)}
Certain activity sets are designated as clusters (e.g., shared lectures across groups, cross-major clusters).
For each cluster $C$ with leader $\ell$, enforce equal start times:
\begin{align}
\mathrm{start}(a) = \mathrm{start}(\ell) \qquad \forall a\in C.
\end{align}

\subsection{Soft Objective (Optional)}
When the balanced objective mode is enabled, the model minimizes
\begin{align}
\min_{x}\; F(x)
&=
w_{\mathrm{free}}\,P_{\mathrm{free}}(x)
+w_{\mathrm{gap}}\,P_{\mathrm{gap}}(x)\nonumber\\
&\quad
+w_{\mathrm{late}}\,P_{\mathrm{late}}(x)
+w_{\mathrm{stab}}\,P_{\mathrm{stab}}(x)\nonumber\\
&\quad
+w_{\mathrm{room}}\,P_{\mathrm{room}}(x)
+w_{\mathrm{same}}\,P_{\mathrm{same}}(x),
\label{eq:softobj}
\end{align}
where each \(P_{\bullet}(x)\) is a nonnegative linearized penalty term:
free-day shortfalls, intra-day gaps, late starts, week-to-week instability, and
room inconsistency (when rooming is integrated in CP), plus weekly concentration
of same-course LEC/TUT sessions for each group.
All weights \(w_{\bullet}\ge 0\) are configurable at the instance level.

For the concentration component, let
\begin{align*}
N_{g,c,k,w}=\#\{a\in\mathcal{A}:{}&g\in\mathcal{G}(a),\\
&\mathrm{course}(a)=c,\ \mathrm{kind}(a)=k,\\
&\mathrm{week}(a)=w\}.
\end{align*}
for \(k\in\{\mathrm{LEC},\mathrm{TUT}\}\), and define
\[
P_{\mathrm{same}}=\sum_{g,c,k,w}\max(0,N_{g,c,k,w}-1).
\]
In the CP model, weeks are fixed per activity, so this term is constant.
The implementation records the constant to keep CP and exact post-solve score scales aligned; it becomes decision-relevant only in editing procedures that permit week changes.

\paragraph{Fairness-first profile.}
Let $B_g(x)$ be the weighted group-specific portion of the soft penalty and let $U$ be a proven upper bound on all secondary penalties.
The fairness profile minimizes
\begin{align}
(U+1)\max_{g\in\mathcal{G}}B_g(x)+F(x).
\end{align}
Because the multiplier exceeds the maximum possible secondary change, this scalarization exactly implements the lexicographic order ``minimum worst group burden, then minimum total penalty''.
Reported fairness diagnostics also include the 90th percentile and total burden, Gini inequality, Jain equality, and burden per scheduled slot.

If objective mode is disabled, the solver runs in pure feasibility mode:
find any \(x\) satisfying all hard constraints.


\section{Certificate-Guided Room Decomposition}
The implementation retains integrated CP rooming and greedy post-processing as baselines, and adds \texttt{room\_mode=decomposed} as the research method.
Strict CP rooming now uses every eligible room by default; candidate truncation is explicit and opt-in.
For a co-located cluster, capacity is computed over the union of participating groups rather than per activity.

\subsection{Fixed-Time Room Jobs}
Given a master schedule, co-located activities are collapsed into room jobs $j\in\mathcal{J}$.
Each job has a fixed week, day, start, duration, union robust-demand requirement $D_j$, and candidate domain $R_j$.
The domain is the intersection of member eligibility after room type, specialization, capacity (when hard), availability, closures, and room locks are applied.
Binary $y_{j,r}$ selects exactly one candidate room:
\begin{align}
\sum_{r\in R_j}y_{j,r}=1 \qquad \forall j\in\mathcal{J}.
\end{align}
Per-room/per-slot at-most-one constraints enforce physical occupancy.
Additional pairwise constraints enforce travel buffers for jobs sharing staff or groups, and equality constraints implement an optional repeated weekly room pattern.
The room objective prioritizes tutorial room type, capacity overflow when capacity is soft, capacity waste, and deterministic room identifiers.
In objective-free runs the exact subproblem stops at feasibility rather than spending the remaining budget proving room-waste optimality; objective-enabled runs retain the room objective.

paragraph{Support-aware feedback proxy.}
For a fixed incumbent-majority room $r_c$ for each course $c$, the compatibility feedback scalar counts event-level collisions at $(q,r_c)$.
The opt-in ablation instead resolves every period--majority-room collision deterministically, retains one capacity-prioritized event, and counts each displaced course once over the timetable.
This distinct-course fragmentation count is only a candidate-ordering proxy: it is neither the official room-stability objective nor an exact fixed-time room solve.
All resulting candidates are lifted to physical rooms and accepted only after full hard validation and canonical official rescoring.

\subsection{Hall-Deficiency Certificates}
For occupied slot $q=(w,d,s)$, form a bipartite graph between active jobs $\mathcal{J}_q$ and their eligible rooms.
A maximum matching is computed.
If it does not cover every active job, alternating reachability from unmatched jobs yields a witness $J'\subseteq\mathcal{J}_q$ with neighborhood $N(J')$ satisfying
\begin{align}
|J'|>|N(J')|.
\end{align}
By Hall's theorem, those jobs cannot receive distinct rooms at $q$.
The certificate stores the activities, representative cluster activities, candidate rooms, slot, deficiency, active demand mode, required capacities, and a counting explanation.

\subsection{Master Cut}
Let $D(j,t)$ be the exact candidate-room domain of representative job $j$ when started at $t$, after cluster membership, aggregate demand, room type, specialization, room locks, interval availability, and closures are applied.
Retain the fixed-time certificate's witnessed room set $R^{\mathrm w}$ as evidence, and define the option-expanded literal set
\begin{align}
U(q,R^{\mathrm w})=
\{(j,t):q\in I(j,t),\;D(j,t)\subseteq R^{\mathrm w}\}.
\end{align}
Only starts whose complete candidate domain is contained in the witnessed room set are counted.
Their actual neighborhood is reconstructed separately as
\begin{align}
\Gamma(U)=\bigcup_{(j,t)\in U(q,R^{\mathrm w})}D(j,t).
\end{align}
The master receives the standard option-expanded Hall inequality
\begin{align}
\sum_{(j,t)\in U(q,R^{\mathrm w})}x_{j,t}
\le |\Gamma(U)|. \label{eq:contextual-hall}
\end{align}
The witness $R^{\mathrm w}$ therefore limits which alternatives can enter the derivation, while the right-hand side is the stronger reconstructed neighborhood cardinality, not $|R^{\mathrm w}|$.
For any selected subset of those alternatives, distinct active jobs require distinct rooms drawn from a subset of $\Gamma(U)$, which proves~\eqref{eq:contextual-hall}; the master's at-most-one start constraint prevents two alternatives of one job from being selected together.
The incumbent witnessed starts must all satisfy the subset rule, and the contextual inequality is emitted only when it excludes the incumbent and is stronger than its exact incumbent-start nogood.
Non-incumbent starts whose domain expands outside $R^{\mathrm w}$ are excluded from $U$.
If an incumbent domain is not contained in $R^{\mathrm w}$, if any domain proof is uncertain, or if infeasibility arises only after travel or repeated-room coupling, the solver emits the universally safe exact incumbent-start nogood instead.
Each cut records its kind, witness slot and rooms, derived $\Gamma(U)$, right-hand side, included and excluded start counts, proof rule, and fallback reason.
The master is re-solved and the exact room subproblem is repeated until a full assignment is found, the master proves infeasible, or the shared wall-clock/round limit is reached.

\subsection{Replay and Integrity Lineage}
Before adding a contextual Hall inequality, a separate checker replays the serialized certificate and derivation without consulting the constructed CP model.
It verifies the certificate schema and deficiency, fingerprints the room-domain input context, reconstructs every effective room domain from the stated members, time, duration, and master start domains, checks that the counted option set is complete, recomputes $\Gamma(U)$ and the right-hand side, and confirms that the inequality excludes the incumbent.
Malformed, incomplete, or inconsistent evidence is never inserted as a contextual cut; when the representative incumbent starts are available, the cut layer uses the exact-nogood fallback.
The replay path is independent of the cut-construction code, although both deliberately call the same authoritative effective-domain routine; it is not a second implementation of room eligibility.

Canonical JSON payloads receive content-addressed SHA-256 identifiers for the certificate, room context, inequality, and complete derivation.
The certificate-derived LNS signal retains these identifiers, the witness and derived room sets, and the domain assumptions; a selected certificate neighborhood records the source certificate, cut, and derivation identifiers in its trace.
These hashes provide deterministic integrity and lineage, not authorship or authenticity: they are neither digital signatures nor protection against an actor replacing an artifact and all of its reference hashes.

\subsection{Explanation and Repair Reuse}
The Hall certificate supports the statement ``$k$ simultaneous room jobs can use only $m<k$ rooms'' and recommends moving at least one certified job outside the slot.
For repair, certified activities are expanded through a typed activity hypergraph linking shared groups, staff, courses, rooms, clusters, distribution constraints, and temporal partition boundaries.
Activities outside the selected neighborhood are fixed through CP-SAT assumptions on a reusable exact model; an infeasible assumption core is converted into another typed search signal.
A discounted-UCB controller chooses between certificate, penalty, partition-boundary, and random families at several target sizes.
Only schedules accepted by the separate repository-local post-solve validator and found to be strictly improving replace the incumbent, and every exact-repair round records build, assumption-setup, search, validation, objective, bound, and local-optimality evidence.
Cut-guided LNS, explanation-guided repair, adaptive LNS, and bandit selection have prior art; the benefit of this integrity-linked systems integration remains an experimental question pending the prescribed ablations.


\section{Local Search Improvement Heuristic}
After a feasible timetable is produced by CP-SAT, an optional local-search improver attempts to reduce a composite soft penalty.
This follows the common pattern of feasibility-first construction plus neighborhood-based improvement widely used in course timetabling~\cite{lewis2008metaheuristics,bettinelli2015overview}.

\subsection{Neighborhood and Move Generation}
The improver is two-stage. First, it runs a deterministic group/week sweep to exploit obvious structural gains. Second, it runs stochastic local search for additional refinements.

\paragraph{Deterministic group/week sweep.}
For each week and for high-penalty groups first, the improver searches feasible relocations and applies the best improving move found.
Besides single-cluster relocations, it evaluates bounded compound relocations (default: up to 2 chained time moves), so a small temporary loss in step 1 can still be accepted when the total combined delta is better.
This pre-pass is explicitly bounded by move count, pass count (default: up to 3 week-wise passes), and wall-clock budget to preserve responsiveness.

\paragraph{Stochastic stage.}
After the deterministic pass, each iteration selects high-pressure activities (or co-start clusters) and evaluates one of:
\begin{itemize}[leftmargin=*]
  \item \textbf{Time move}: change $(d,s)$ within the same week, respecting hard constraints checked by incremental occupancy tables.
  \item \textbf{Room move}: change room assignment to another eligible room (when room assignments exist), respecting room eligibility and availability.
\end{itemize}

\noindent
Compared to purely random local search, candidate generation is guided:
\begin{itemize}[leftmargin=*]
  \item groups and staff are scored by current soft-penalty ``badness''; move selection prioritizes activities attached to the worst entities first,
  \item before stochastic sampling, a deterministic week-wise sweep targets worst group-weeks directly,
  \item activities are further ranked by a pressure score (gaps, thin days, late starts, excess active days, staff day pressure),
  \item for a selected activity/cluster, the search scans feasible candidates and keeps the best local delta found,
  \item room moves use intersection of allowed-room sets for clustered activities to avoid invalid proposals early.
\end{itemize}

\subsection{Feasibility Preservation}
Candidate moves are filtered by:
\begin{itemize}[leftmargin=*]
  \item group and staff non-overlap in the affected slots,
  \item staff availability, and optional load caps,
  \item room non-overlap and eligibility (if a room is assigned),
  \item lock constraints (locked activities cannot be moved).
\end{itemize}

\subsection{Fast Delta Evaluation}
To keep improvement interactive, the improver avoids full rescoring on every candidate.
It evaluates local deltas on affected components only:
\begin{itemize}[leftmargin=*]
  \item group-week terms (free-day deficits, gaps, thin/single days, late starts, active-day penalty),
  \item staff-week free-day term,
  \item adjacent week-pair stability term,
  \item room-consistency keys touched by the move.
\end{itemize}
An exact global penalty recomputation is triggered periodically for numerical robustness.
Entity badness scores (groups/staff) are refreshed periodically as the schedule changes, so priority targeting remains aligned with the current worst parts of the timetable.

\subsection{Acceptance Rule}
We use a simulated annealing style rule.
Let $\Delta$ be the change in penalty.
Moves with $\Delta\le 0$ are accepted; moves with $\Delta>0$ are accepted with probability $\exp(-\Delta/T)$ where $T$ follows an exponential cooling schedule from $T_0$ to $T_1$.

\subsection{Restart and Diversification Strategy}
To reduce plateauing, the improver uses a bounded multi-start policy inside the same time budget:
\begin{itemize}[leftmargin=*]
  \item if no improvement is observed for a restart window, the search restarts from the current elite solution,
  \item a short diversification ``kick'' applies a few random feasible perturbations,
  \item guided search then resumes from this perturbed elite state.
\end{itemize}
This preserves hard-feasibility checks while improving search coverage under strict runtime limits.


\section{System Design and Implementation}
The system follows a feasibility-first decomposition that separates three concerns:
time assignment, room assignment, and quality improvement.

\subsection{Decomposition Strategy}
\begin{itemize}[leftmargin=*]
  \item \textbf{Stage 1 (time feasibility):} assign each activity to a valid week/day/slot while satisfying all hard constraints.
  \item \textbf{Stage 2 (rooming):} co-optimize rooms inside CP (\texttt{cp\_rooms}), assign rooms afterward (\texttt{greedy}), or iterate between an exact room subproblem and certificate cuts (\texttt{decomposed}).
  \item \textbf{Stage 3 (quality):} improve a feasible schedule using local search under a soft-penalty objective.
\end{itemize}
This decomposition provides a practical trade-off between strictness and runtime.

\subsection{Adaptive Week-Partitioned Scale Mode}
For large instances, the operational \texttt{partitioned} mode exploits the fact that group, staff, room, and generic-resource conflicts are indexed by teaching week in the current model.
It first runs global structural checks for course totals, staff competency, and block-teaching invariants.
Before partitioning, every required distribution relation is inspected: week-separable aggregate and non-overlap rules are split exactly, fixed week relations are decided statically, and a relation that genuinely couples decisions across weeks causes a fail-fast error, as does a forced repeated-weekly pattern.
Cross-week precedence is accepted only when its fixed week order already satisfies the rule; reverse order is rejected.

Each remaining week is compiled into a compact objective-free CP master.
Up to four week partitions run concurrently by default, each with one CP-SAT search worker to bound memory multiplication.
A valid fast room assignment completes the partition.
If room extraction fails, only that week is rebuilt with the exact certificate-guided room subproblem and the remaining shared wall budget.
Required same-room or different-room relations select the exact inner mode immediately.
All rows are merged only when every partition is feasible, after which the original full instance validator rechecks completeness and every hard rule.
The mode does not report a global soft-objective bound: stability, fairness, and other cross-week quality terms require monolithic optimization or a separately reported improvement phase.

\subsection{Reliability Policy}
To reduce solver timeouts in difficult instances, the runtime applies a bounded retry policy:
\begin{enumerate}[leftmargin=*]
  \item objective-enabled CP attempt (quality-aware),
  \item objective-disabled retry in the same room mode (feasibility-oriented),
  \item optional fallback from strict rooming to greedy rooming.
\end{enumerate}
The operational policy remains available, but research comparisons disable it and analyze primary attempts only.
One-worker configurations are intended to be reproducible for fixed seeds; multi-worker CP-SAT runs are labeled separately.

\subsection{Budgeted Anytime Mode}
For interactive use, total solve time budget $B$ is split into
a feasibility phase $B_f$ and an improvement phase $B_i$:
\begin{align*}
B_f + B_i &= B,\\
B_i &\approx 0.2B\ \text{(capped)},\\
B_f &\approx 0.8B\ \text{(floored)}.
\end{align*}
The first phase prioritizes obtaining a valid schedule quickly; the second phase runs short local-search slices to improve penalty.
Each improvement slice uses a capped deterministic pre-pass (group/week sweep with bounded compound lookahead) followed by guided stochastic neighborhoods with bounded restarts/diversification when the search plateaus.

\subsection{Operational Parameters}
The main runtime parameters are:
\begin{itemize}[leftmargin=*]
  \item room mode (\texttt{cp\_rooms}, \texttt{decomposed}, \texttt{partitioned}, or \texttt{greedy}),
  \item objective on/off,
  \item time limit and worker count,
  \item local-search iteration/time budget.
\end{itemize}
The same solver parameterization is used across batch and interactive runs, supporting consistent interpretation of results.

\subsection{Interactive Administration Layer}
The desktop UI exposes model-consistent administrative edits:
\begin{itemize}[leftmargin=*]
  \item hold-and-move operations across days/slots and also across weeks (subject to hard checks),
  \item immediate conflict diagnostics for candidate targets before committing moves,
  \item real-time recomputation of global soft penalty deltas for candidate placements.
\end{itemize}
Staff week-availability checks are enforced in the same validator path used by solver outputs.

\subsection{Custom Instance Authoring}
The custom generator tab supports:
\begin{itemize}[leftmargin=*]
  \item per-program group/course overrides,
  \item per-course LEC/TUT/LAB count overrides (including lec-only, tut-only, and lab-only patterns inferred directly from counts),
  \item staff day/week availability and teaching-course mapping,
  \item room capacity policy as either categorical or exact numeric mode.
\end{itemize}
Custom generation configurations can be saved locally for reuse and exported/imported as JSON profiles.

\subsection{Generation-Time Solvability Safeguard}
In custom scenarios, users may define room sets that accidentally leave some activities without any eligible room
(because of type, specialization tag, or capacity).
To avoid trivial unsatisfiable instances from such configuration mismatches, generation applies a coverage safeguard:
if an activity has zero eligible rooms, a minimal fallback room of the required class is injected automatically.
This does not remove hard constraints during solving; it only ensures that each generated activity has at least one feasible room candidate.


\section{Experimental Evaluation}
This version defines the evaluation protocol and evidence gates.
It reports a generated feasibility-architecture study, an official-validator-checked ITC-2007 breadth comparison, a single-seed fixed-time room-dive ablation, and bounded engineering-smoke evidence; the external quality result is negative and the remaining ablations stay open.

\subsection{Quarantined Historical Evidence}
The previously generated JSONL tables are retained only for provenance and are not evidence for this paper.
Two defects invalidate the intended interpretation:
\begin{itemize}[leftmargin=*]
  \item seed labels were not passed into named generators, so seeds 1, 2, and 3 produced identical instances within each mode;
  \item a requested \texttt{cp\_rooms} runtime could include an objective attempt, an objective-free retry, and a greedy fallback, while the requested \texttt{greedy} condition could follow a different path.
\end{itemize}
The old claims of 24/24 reliability, strict-vs-greedy speed ratios, and local-search gains are therefore withdrawn.
The artifact records this status in \path{paper/results_status.json}.

\subsection{Research Questions}
The regenerated study will address:
\begin{itemize}[leftmargin=*]
  \item \textbf{RQ1:} Does certificate-guided decomposition improve feasible-solution rate or time-to-feasible relative to full integrated room CP under an identical total budget?
  \item \textbf{RQ2:} How many master iterations and Hall or generic cuts are required, and how large are the certificates?
  \item \textbf{RQ3:} Do certificate-guided repair neighborhoods reduce repair time and collateral activity changes relative to generic resource-linked and random neighborhoods?
  \item \textbf{RQ4:} What price in total penalty and runtime is paid for lexicographically reducing worst-group burden?
  \item \textbf{RQ5:} Which findings transfer from generated workloads to official-validator-checked ITC-2007 instances and, once the corpus and official validator are available, locally parsed ITC-2019 instances with fixed placements?
  \item \textbf{RQ6:} When week independence is established by preflight, how do adaptive parallel partitions change time-to-feasible and peak memory relative to a monolithic objective-free master?
\end{itemize}

\subsection{Conditions and Controls}
Named generators receive an explicit instance seed.
The same integer is recorded separately as instance, CP-SAT, and local-search seed, and every normalized instance and extracted schedule receives a SHA-256 fingerprint.
Canonical research runs use one CP-SAT worker; multi-worker runs are labeled non-deterministic configurations.

The primary comparison disables retries and fallbacks.
Each algorithm receives the same total wall-clock budget and is analyzed from its primary attempt only:
\begin{itemize}[leftmargin=*]
  \item integrated full-domain \texttt{cp\_rooms};
  \item certificate-guided \texttt{decomposed};
  \item time-master plus \texttt{greedy}, reported as a practical baseline whose post-processing can fail.
\end{itemize}
Portfolio retries are evaluated separately as an operational policy, never merged into an algorithm runtime.

At least 30 effective instance seeds per generated condition are required before inferential claims.
Each row records status, objective, best bound, relative gap, branches, conflicts, solver-reported and external wall time, environment, Git revision and dirty state, decomposition rounds, certificates, and final validation.

\subsection{Statistics}
Feasibility is reported with a binomial confidence interval.
Runtime and penalty use medians and interquartile ranges plus bootstrap confidence intervals.
Algorithms are paired only when their \texttt{instance\_sha256} values match.
Paired differences are accompanied by bootstrap intervals and an exact sign test; effect sizes and raw rows remain available.
Timeouts are not discarded: time-to-feasible analysis treats them as right-censored observations or reports conservative budget-capped values.
No claim is based solely on a $p$-value.

\subsection{Generated Feasibility Gate}
We executed 50 paired \texttt{small\_demo} instance seeds for each room architecture, for 150 primary runs in total, with one worker, a 10~s limit, objective disabled, and no retry, fallback, local search, or adaptive repair.
An effective observation requires a unique normalized-instance SHA-256, a primary feasible result, no fallback, and zero errors from a separate repository-local post-solve validator.
The initial 30 seeds produced only 21--22 effective observations per condition and therefore failed the preregistered gate; rather than weaken the criterion, we preserved that shard and added seeds 31--50.
The combined checked-in, content-hashed shards contain 37--38 effective instances per condition and pass the minimum-30 gate.

\begin{table*}[t]
\centering
\footnotesize
\begin{tabular}{lrrrr}
\toprule
Room architecture & Effective / attempted & Feasible rate (Wilson 95\% CI) & Effective median (s) & IQR (s) \\
\midrule
Integrated \texttt{cp\_rooms} & 38 / 50 & 0.760 [0.626, 0.857] & 3.478 & [2.163, 4.215] \\
Exact \texttt{decomposed} & 37 / 50 & 0.740 [0.604, 0.841] & 1.124 & [0.815, 1.489] \\
Time master + \texttt{greedy} & 38 / 50 & 0.760 [0.626, 0.857] & 0.698 & [0.453, 0.886] \\
\bottomrule
\end{tabular}
\caption{Primary CP search time on feasible generated instances accepted by a separate repository-local post-solve validator. Infeasibility proofs and the single decomposed timeout are retained in the raw rows but excluded from the effective-time summaries.}\label{tab:generated-feasibility}
\end{table*}

On the 37 instances effective for both exact formulations, integrated CP minus decomposition had a paired median of 2.177~s, a bootstrap mean-difference 95\% interval of [1.652, 2.687]~s, and exact sign-test $p=5.53\times10^{-10}$.
On the 37 exact-decomposition/greedy pairs, decomposition minus greedy had a paired median of 0.405~s and a bootstrap interval of [0.394, 0.776]~s.
Thus the exact room decomposition materially reduces time-to-validated-feasible relative to the integrated room model on this generated condition, while greedy post-processing remains faster.
The feasibility rates do not establish a reliability advantage, and objective-free soft penalties are exploratory only: the three formulations do not expose commensurate optimization variables or bounds under this protocol.
The two raw shards and the combined analysis are content-hashed in \path{paper/evidence/generated_feasibility_50x3_2026-08-11_analysis.json}. % chktex 29
Each raw row records a base Git revision and \texttt{git\_dirty=true}, but the runner did not capture a file-by-file manifest of the dirty working tree.
The checked-in rows and shard hashes therefore preserve result integrity and configuration provenance, not an exact reconstruction of the executed source snapshot.

\subsection{Certificate and Repair Evaluation}
For each infeasible room subproblem we record certificate type, $|J'|$, $|N(J')|$, deficiency, cut strength, and whether the next master solution removes the witnessed overload.
Certificate soundness is checked by a separate replay path that reconstructs the filtered bipartite graph and verifies $|J'|>|N(J')|$.

Disruption trials sample used staff-week and room-week pairs with a seeded protocol.
Outcomes include affected and unresolved activities, recovery rate, hard conflicts, repair wall time, and number of activities changed from baseline.
The implemented script currently evaluates deterministic staff/room reassignment baselines; a publication claim about certificate-guided repair requires the planned controlled neighborhood ablation.

\subsection{External Validity}
The ITC-2007 importer preserves curricula as zero-size conflict groups, adds a separate enrollment group for capacity, and maps course-specific unavailability to activity-specific forbidden starts.
Room capacity is soft in the imported variant.
The external scorer implements the official ITC-2007 weights: one per excess student, five per missing working day, two per isolated curriculum lecture, and one per additional room used by a course~\cite{digaspero2007itc}.

\paragraph{Official ITC-2007 breadth.}
Under the recorded one-worker, 10~s, seed-17 protocol, two fresh source-identical passes produced 84 officially feasible, zero-hard-violation schedules across all 21 competition instances and both solvers, with exact Planora internal/official component agreement.
Pass A yields Planora/CPSolver/tie counts 9/10/2 and mean objectives 156.95/173.86; pass B yields 14/7/0 and means 158.14/178.38~\cite{muller2009hybrid,unitimeitc2007}.
Only 14 of 21 instance winners remain unchanged across the two passes, so the lower Planora mean in both passes is promising breadth evidence but not a deterministic or general superiority result.
The consolidated source/tool/validator-bound record is \path{paper/evidence/itc2007_fresh_redundant_matrix_2026-08-13.json} (SHA-256 \path{ef02acc47d02c59462a00a4ba57b25b1e39ce02a6d98dd7061c1e28cc269409e}).
The older 0/21 breadth artifact remains historical evidence and is not used as the current quality conclusion.

\paragraph{Fixed-time room-dive ablation.}
Fixing all event starts and reoptimizing rooms is directly related to established timetabling dives and exact room matching~\cite{burke2010diving,lu2008hybrid}.
We therefore test it as an engineering ablation, not a novelty claim.
An earlier breadth attempt is retained as a failed gate: the ON condition was better on five instances, tied on 14, and worse on two, with reported sums 19,711 ON versus 20,360 OFF, but seven of its 21 ON runs exceeded the strict total deadline; those rows are excluded from the corrected result.
After adding explicit dive and outer-completion reserves, a fresh counterbalanced matrix with source SHA-256 \path{a63732de42d9e0053ea196f46a410024b85c006e4cc6d7f91e5529f3bd1294cb} produced 42/42 officially feasible, zero-hard-violation schedules, with 42/42 internal/official component agreement, source stability, and zero strict total-deadline overrun.
Here, source stability means that every matrix record matched that frozen source snapshot; the snapshot predates the later repeat-week greedy-to-decomposed correctness-only change, so the matrix is immutable historical evidence rather than evidence for the exact current checkout.
ON was better on three pairs, tied on 18, and worse on none; its aggregate objective was 19,870 versus 20,457 OFF, a 587-point (2.869\%) reduction.
However, 13 dives were attempted, only one was accepted, 12 returned no accepted improvement, and eight were skipped for insufficient reserved budget.
Only \texttt{comp01} records a direct within-run accepted-dive improvement (499 to 253, $-246$); the other two paired differences arose from different initial incumbents and are not attributed to the dive.
Against the immutable CPSolver rows, the ON condition still loses all 21 instances and is 5.541 times the CPSolver aggregate (19,870 versus 3,586).
This is a single-seed ablation with no competitor superiority, so the feature remains disabled by default.
The final evidence report is \path{reports/itc2007_fixed_time_room_dive_breadth_final_v2_seed17_2026-08-11.json} (SHA-256 \path{24fccbf81d78a1371d625f42c94e1ff1b7d8e841633874ddb788a552d7146be6}), backed by \path{output/itc2007-room-dive-breadth-seed17-counterbalanced-final-v2/matrix_index.json} (SHA-256 \path{aba9413209699d2eff0db431f90e971dea3dada0e675420f3d9cefeb6bc8f2f4}).
The preceding negative report is preserved at \path{reports/itc2007_fixed_time_room_dive_breadth_seed17_2026-08-11.json} (SHA-256 \path{43560980d667a5495dda0b597ae4942f4e355ed94ec041dbdb0fcd3cd2d912d7}), and the intervening targeted deadline-correction gate is \path{reports/itc2007_fixed_time_room_dive_deadline_gate_v2_2026-08-11.json} (SHA-256 \path{d595fb73473c804e66eceba47301cc12fb4c7522ee17e1afb75c578bd624fcc2}).

\paragraph{Distinct-course room-support mechanism check.}
The opt-in support proxy was evaluated on the retained Planora-native \texttt{comp10} incumbent, not as a new cold solve.
Two fixed-checkpoint, one-worker calls under a 2.2~s absolute post-incumbent deadline produced byte-identical official-valid schedules: 109 $(0,20,48,41)$ became 96 $(0,20,46,30)$ in 1.736~s and 1.814~s, with zero deadline overrun and exact internal/official component agreement.
The retained CPSolver artifact independently scores 82, so this is a 13-point mechanism gain but not a competitor win.
Increasing only the collision scalar weight produced 102, which is worse than the distinct-course proxy and is retained as a negative ablation.
The evidence record is \path{paper/evidence/itc2007_support_proxy_comp10_2026-08-13.json}; timing excludes construction of the retained incumbent and cannot support an end-to-end runtime claim.

\paragraph{Selected post-incumbent quality checks.}
After the breadth diagnosis, five representation-derived portfolios were tested on retained Planora-native incumbents, with strict independent validation after every accepted handoff.
They produced lower scores than the retained independently validated CPSolver artifacts on \texttt{comp02} (133 versus 137), \texttt{comp03} (162 versus 167), \texttt{comp19} (154 versus 156), \texttt{comp21} (193 versus 194), and Exam set~1 (5790 versus 5824).
The bounded Planora calls took 4.076, 1.947, 2.240, 2.143, and 4.948 seconds respectively, all with zero reported deadline overrun.
These are selected post-incumbent mechanism results: incumbent construction is excluded, comparator runtime is not matched, and the cases were selected after diagnosis.
Moreover, the \texttt{comp19} and \texttt{comp21} records bind earlier implementation hashes that predate a later default-preserving API extension.
They therefore demonstrate concrete basin-repair capability, not end-to-end or corpus-wide superiority.
The complete score, timing, artifact-hash, and source-match boundary is recorded in \path{paper/evidence/selected_post_incumbent_quality_2026-08-13.json}.

\paragraph{ITC-2019 conversion and bounded competitor matrix.}
The ITC-2019 component parses the official XML structure and supports native Cartesian and factorized placement/sectioning formulations, official-format solution interchange, and an independent repository-local validator/scorer~\cite{itc2019official,itc2019format,muller2025itc2019}.
A same-host matrix executed the exact 30 competition instances shown on the authenticated Results page against Planora, Gashi simulated annealing, UniTime/CPSolver, and the Lemos MaxSAT finalist: one seed, one repetition, one worker, one CPU affinity, and a nominal 10~s solver budget, for 120/120 recorded conditions.
Only three conditions emitted complete locally valid outputs, all from CPSolver: 18,390 on \texttt{mary-spr17}, 282 on \texttt{lums-spr18}, and 35,608 on \texttt{nbi-spr18}; all are worse than the authenticated competition best scores of 14,910, 95, and 18,014.
Authenticated official-validator uploads agreed with the independent total and all four components for 3/3 outputs.
A separate four-solver \texttt{lums-sum17} smoke produced score 4 for every solver and official agreement for 4/4 uploads.
The tools expose different timeout and memory mechanisms, including a CPSolver-only 768~MiB Java heap cap in this executed matrix, so this is descriptive bounded quality evidence, not an equal-wall runtime, speed, or general superiority result.
The matrix and claim boundary are recorded in \path{output/itc2019-open-source-competition-30-seed17-10s-20260813/report.json} and \path{docs/ITC2019_COMPETITOR_MATRIX_2026-08-14.md}.

\subsection{End-to-End Performance Measurement}
The shared-backend runner exercises generation, JSON serialization, session creation, solving, independent validation and scoring, improvement, final validation/scoring, and response serialization.
On the 3,888-activity proxy, five measured runs after one warmup were valid in each transport. Median total time was 17.668~s embedded, 28.080~s through a freshly started HTTP API process, and 6.234~s through a persistent HTTP service; the latter benefits from the exact shared solve cache.
The embedded medians were 12.105~s for solve and 4.039~s for Improve. Every run returned all 3,888 activities with zero hard conflicts and zero validation errors.
Replacing twelve full-instance deep copies with isolated shallow partition views reduced observed partition model construction from about 2.00~s to a 0.313~s median (approximately $6.4\times$); harder weeks now enter the four-worker queue first.
These are single-host engineering measurements, not latency guarantees. The methods and raw evidence paths are recorded in \path{docs/PERFORMANCE_3888_2026-08-14.md}.

\paragraph{GIU-scale path.}
The local synthetic target-case generator, configured with the historically derived \texttt{giu\_target} policy layer, contains 3,888 activities, 12 weeks, 29 groups, 67 synthetic staff records, and 35 generated rooms.
These resource counts are scale-test inputs, not institutional facts.
In particular, the historical SS23 source contains 15 distinct nonblank room labels and 146 rows whose room is missing; the importer preserves those missing values as row-scoped placeholders, so its 161 modeled room records must not be interpreted as 161 physical rooms.
The synthetic scale instance compiles to 3,888 compact start integers over 105,510 admissible domain values.
The current historical replay reports 148 validation errors: 93 synthetic-staff overlaps, 39 inferred room-type or specialization mismatches, and 16 apparent room overlaps.
Consequently, the hash-locked calibration in \path{paper/evidence/giu_ss23_calibration.json} (SHA-256 \path{c2620b38ed4535fbbeb55ca66ec3c031cf346a1a5dc6c66d62cc545c5df53b73}) establishes historical source provenance and a conservative research preset only; historical schedule/model agreement and current institutional calibration remain open.
The repeated performance evidence establishes the implementation's current synthetic scale boundary, not institutional representativeness or an algorithmic generalization claim.

\newpage
\subsection{Current Evidence Status}
\begin{center}
\centering
\scriptsize
\renewcommand{\arraystretch}{0.92}
\begin{tabular}{p{0.42\columnwidth}p{0.46\columnwidth}}
\toprule
Evidence item & Status \\
\midrule
Effective seeded generation & Implemented and unit-tested \\
Primary-attempt isolation & Implemented \\
Bounds, gaps, hashes, environment & Implemented \\
Paired bootstrap/sign-test analysis & Implemented and unit-tested \\
Exact room subproblem, contextual Hall cuts, and safe nogood fallback & Implemented, integration-tested, and exhaustively checked on tiny room domains \\
Certificate explanation/repair scope & Implemented and unit-tested \\
ITC-2007 variant ingestion and official-validator bridge & 21/21 officially feasible with internal/official component parity in the source-hashed breadth run \\
ITC-2007 competitor quality & CPSolver wins 21/21; Planora/CPSolver aggregate $=5.625$ in the rescue comparison \\
Fixed-time room-dive ablation & Corrected single-seed matrix: 42/42 official-valid and deadline-compliant; one accepted direct gain; default off \\
Distinct-course room-support proxy & Default-off mechanism evidence: two byte-identical official-valid \texttt{comp10} improvements, 109 to 96; retained CPSolver is 82 \\
Selected post-incumbent checks & Lower than retained CPSolver on four CTT cases and Exam set~1; diagnostic selection and non-equal timing prohibit a general superiority claim \\
Portable institutional policy/demand metrics & Implemented and unit-tested \\
Adaptive week partitioning and safeguards & Implemented, unit-tested, and locally smoke-tested \\
ITC-2019 native path and competitor matrix & 120/120 bounded conditions recorded; 3 complete valid outputs, with authenticated official agreement 3/3; effective breadth remains open \\
3,888-activity end-to-end performance & Five valid runs per transport; embedded median 17.668~s and persistent-HTTP median 6.234~s on one host \\
Generated feasibility-architecture study & 50 paired attempts; 37--38 unique effective instances per condition; generated-feasibility gate passed \\
Quality and certificate-guided neighborhood ablations & Not yet publication-ready \\
Historical GIU calibration & Source-hash locked; 148 replay errors and current institutional sign-off remain open \\
\bottomrule
\end{tabular}
\refstepcounter{table}\label{tab:evidence-status}
\par\smallskip
\footnotesize\textbf{Table~\thetable:} Separation of implemented infrastructure from empirical evidence.
\end{center}

\FloatBarrier

\section{Limitations and Future Work}
The most important limitations are methodological rather than cosmetic:
\begin{itemize}[leftmargin=*]
  \item \textbf{Broad novelty is falsified.} Decomposition, arbitrary-bipartite Hall inequalities, Benders-guided LNS, explanation-guided exact repair, adaptive neighborhoods, and fairness have direct prior art. The remaining claim is a systems-integration hypothesis, not a new mathematical inequality or metaheuristic class.
  \item \textbf{Generated feasibility evidence is narrow.} Historical rows remain invalid. The replacement 50-instance study passes its minimum-30 effective-instance gate and supports a generated feasibility-latency comparison, but it does not establish objective quality, certificate-guided repair benefit, institutional transfer, or superiority over external solvers. Its raw rows record a base Git revision and dirty-tree flag but no file-level manifest of the executed dirty source, so the content-hashed result shards are not a complete source reconstruction.
  \item \textbf{Hall certificates are partial explanations.} A contextual cut includes only starts whose exact room domain is proven to be a subset of the witnessed room set; otherwise the solver emits the exact incumbent nogood. Travel, repeated-room, and lock coupling may therefore require a less specific explanation.
  \item \textbf{Fairness depends on burden semantics.} Lexicographic optimization is exact for the encoded score, but the score itself is a normative choice. Gini or Jain values cannot establish stakeholder fairness without elicitation and sensitivity analysis.
  \item \textbf{Fresh external quality is competitive but repeat-sensitive.} Two source-identical, official-validator-backed 42-row ITC-2007 passes close feasibility on all 84 rows. Planora has the lower mean objective in both, but paired counts change from 9/10/2 to 14/7/0 and seven instance winners change. This supports competitive breadth, not deterministic or general superiority; a prespecified multi-seed design remains required.
  \item \textbf{The room-dive evidence is only single-seed and snapshot-specific.} The corrected 42-run ITC-2007 matrix has official-validator/component agreement and zero strict deadline overrun; ON is non-worsening across the 21 pairs and its aggregate is 2.869\% lower. Only one of 13 attempted dives was accepted, however, and only that \texttt{comp01} trace supplies a direct within-run gain. Its source hash predates a later repeat-week greedy-to-decomposed correctness-only change, so it is immutable historical evidence rather than exact-current-checkout evidence. An earlier breadth run exceeded the total deadline on seven ON cases and is retained as a failed gate. Multi-seed confirmation is required, and the feature remains disabled by default.
  \item \textbf{The support proxy is a post-incumbent mechanism check.} Its two \texttt{comp10} runs are deterministic, official-valid, and deadline-compliant, but they share one retained incumbent, instance, and seed. The proxy improves 109 to 96 yet still trails CPSolver's retained 82. It remains opt-in pending a held-out multi-instance comparison against the event-collision baseline; neither the result nor the proxy supports a new-neighborhood, exact-objective, or end-to-end superiority claim.
  \item \textbf{Selected comparator wins are diagnostic, not a paired matrix.} Four CTT incumbents and one Exam incumbent were improved below retained CPSolver scores, but the Planora calls begin after parsing and incumbent construction, use case-specific bounded portfolios selected after diagnosis, and have no equal-runtime comparator execution. Two CTT proofs also bind earlier source hashes after a default-preserving API extension. These observations motivate a preregistered fresh matrix; they do not overturn the negative 21-instance breadth result.
  \item \textbf{ITC-2019 official agreement is sparse rather than absent.} Authenticated uploads agree for all three complete locally valid outputs in the 120-condition competition matrix and for all four outputs in the separate \texttt{lums-sum17} parity smoke. The competition matrix still has only 3/120 effective outputs, uses one seed, and cannot claim equal resources because upstream timeout semantics differ and CPSolver alone had a 768~MiB heap cap. Agreement on emitted artifacts does not establish effective publication-scale coverage.
  \item \textbf{Policy presets require institutional validation.} The GIU target preset has a hash-locked, limitation-aware calibration against a local Berlin Spring 2023 timetable snapshot, but this establishes only historical provenance and observed grid/scale facts. The current projection has 148 replay errors (93 synthetic-staff overlaps, 39 inferred room-type or specialization mismatches, and 16 apparent room overlaps), so even historical schedule/model agreement is open. It is not an official current policy statement, and the institutional questionnaire and sign-off gate remain open.
  \item \textbf{Robust demand depends on forecast quality.} Worst-case, quantile, and budgeted policies are implemented consistently, but scenario generation and service levels remain institution-specific modeling choices.
  \item \textbf{Synthetic transfer risk.} The local SS23-derived preset uses synthetic staff and inferred room semantics and is suitable for scale testing, not institutional ground truth. In particular, the importer's 161 modeled room records comprise 15 nonblank historical labels plus 146 row-scoped placeholders for missing labels; they are not evidence of 161 physical rooms.
  \item \textbf{Repair evidence is incomplete.} Deterministic outage baselines exist, but the certificate-guided neighborhood still needs controlled comparison against generic and random neighborhoods.
  \item \textbf{Partitionability is conditional.} The adaptive scale mode is exact only for the represented feasibility problem after hard cross-week decision coupling has been excluded. It can split week-separable aggregate/non-overlap rules and decide fixed week relations statically, but it rejects repeated weekly patterns and genuinely coupled cross-week relations. It does not produce a global quality bound.
  \item \textbf{Scalability and end-to-end performance.} The repeated 3,888-activity runs establish one SS23-calibrated synthetic proxy on one host: embedded median 17.668~s, fresh-process HTTP median 28.080~s, and persistent-HTTP median 6.234~s. The warm path includes an exact cache hit, the HTTP client RSS excludes the API process, and the proxy is not a direct GIU reconstruction. These measurements are not service-level guarantees or proof of current institutional capacity.
  \item \textbf{Release hardening is not deployed validation.} A repository-local formal security equivalent closed high-severity token, regular-expression, request/SSE, and solver-admission findings, and an authorized Edge pass covered responsive, keyboard, labeling, focus, overflow, reduced-motion, and WCAG-AA contrast behavior. These checks do not test deployed TLS, reverse-proxy policy, production secrets, network isolation, or institutional operations.
\end{itemize}

The remaining publication gates are prespecified multi-seed external quality, authenticated ITC-2019 agreement with effective breadth, GIU institutional sign-off, and complete ablations for cut strength, explanations, fairness, and repair neighborhoods.


\section{Conclusion}
Planora now implements a certificate-guided time/room decomposition in which an exact room subproblem can produce a Hall-deficiency witness, a start-dependent domain-monotone occupancy cut with an explicit proof rule, an administrator-facing explanation, and a bounded repair scope; uncertain cases fall back to an exact incumbent nogood.
The same codebase also supports integrated CP and greedy baselines, adaptive parallel week solving with explicit coupling safeguards and exact fallback, scenario-robust room capacity, portable institutional constraints and policy presets, exact lexicographic worst-group fairness, deterministic research runs, official ITC-2007 soft scoring, paired statistics, disruption trials, and a checksum-verifiable artifact workflow.

The engineering result now includes a defensible generated feasibility-latency study, but it is not yet a completed empirical paper.
Time/room decomposition and fair timetabling are established ideas, and historical project measurements remain quarantined.
The 50-instance generated study meets its effective-instance gate and finds materially lower paired time-to-validated-feasible for exact decomposition than integrated room CP, with greedy rooming faster still; it makes no quality or transfer claim, and its dirty-tree rows lack an exact source-file manifest.
Two fresh official-validator ITC-2007 passes establish feasibility for all 84 rows and competitive breadth: Planora has the lower mean objective in both, while seven changing instance winners and paired counts of 9/10/2 then 14/7/0 reject a deterministic superiority claim.
The corrected fixed-time room-dive matrix is deadline-safe and non-worsening for this single seed, but only one accepted dive has a direct within-run gain; the method remains default-off and established prior art.
The ITC-2019 parser, auto-selected native formulations, separate repository-local validator, and solution interchange are implemented locally. A 120-condition four-solver matrix produced only three complete valid outputs, and authenticated official validation agrees for all three; this closes score-parity authority for those artifacts, not effective breadth or comparator superiority.
The next result must come from the remaining protocol rather than from stronger prose: prespecified multi-seed external quality, certificate and repair ablations, signed institutional review of the 2026 sensitivity scenario, effective ITC-2019 breadth, and independent confirmation of the narrowed systems hypothesis.


\IfFileExists{IEEEtran.bst}{\bibliographystyle{IEEEtran}}{\bibliographystyle{ieeetr}}
\bibliography{refs}

\appendices
\section{Reproducibility Notes}
The project targets Python 3.12 with pinned dependencies in \path{pyproject.toml} and \path{requirements.txt}.

\subsection*{Seeded Primary Experiments}
Canonical runs use one worker and disable retry/fallback mixing:
\begingroup\scriptsize
\begin{verbatim}
.venv/bin/python scripts/run_experiments.py \
  --mode small_demo \
  --room-modes cp_rooms,decomposed,greedy \
  --seeds 1,2,3,4,5,6,7,8,9,10,11,12,13,14,15,\
16,17,18,19,20,21,22,23,24,25,26,27,28,29,30 \
  --workers 1 \
  --use-objective 0 --time-limit 10 \
  --retry-without-objective 0 \
  --cp-rooms-fallback-to-greedy 0 \
  --ls-iters 0 \
  --out paper/evidence/\
generated_feasibility_30x3_2026-08-11.jsonl
.venv/bin/python scripts/run_experiments.py \
  --mode small_demo \
  --room-modes cp_rooms,decomposed,greedy \
  --seeds 31,32,33,34,35,36,37,38,39,40,41,42,\
43,44,45,46,47,48,49,50 --workers 1 \
  --use-objective 0 --time-limit 10 \
  --retry-without-objective 0 \
  --cp-rooms-fallback-to-greedy 0 \
  --ls-iters 0 \
  --out paper/evidence/\
generated_feasibility_supplement_\
31_50_2026-08-11.jsonl
.venv/bin/python scripts/analyze_experiments.py \
  paper/evidence/generated_feasibility_30x3_\
2026-08-11.jsonl \
  paper/evidence/\
generated_feasibility_supplement_\
31_50_2026-08-11.jsonl \
  --minimum-effective-instances 30 \
  --require-publication-ready \
  --out paper/evidence/\
generated_feasibility_50x3_\
2026-08-11_analysis.json
\end{verbatim}
\endgroup
Each row stores separate seeds, instance and schedule SHA-256 values, direct primary-attempt metrics, bounds/gaps, environment, and any decomposition trace.
The checked-in shards have SHA-256 values \path{624370c0545f72ff75a68e0502543dacd22e2975cba1b2f5621559529c50153c} and \path{58bd92b62d4350c3f33d0ac657d0cb07280ab551dc5736f94cb0dab8d6d5a47f}; the combined analysis hash is \path{07ad129999e29d6a39d45e673545d80580f700a3cdecf674ea500ba7a5c718ba}.
The rows record a base Git revision and dirty-tree flag, but not a file-level source manifest; reproducing their exact executed dirty snapshot requires provenance that these artifacts do not contain.

\subsection*{Local Scale Path}
\begingroup\scriptsize
\begin{verbatim}
.venv/bin/python scripts/benchmark_local_app.py \
  --mode giu_target \
  --room-mode partitioned \
  --time-limit 30 --workers 4 \
  --seed 1 --skip-desktop-startup \
  --out giu-partitioned.json
\end{verbatim}
\endgroup
Use \texttt{--warmup-runs} and \texttt{--repeats} for a latency distribution and preserve every raw report.
No repeated raw browser/desktop end-to-end distribution is claimed as checked-in paper evidence at this revision; the existing engineering-smoke file is regression provenance only.

\subsection*{Historical GIU Calibration Evidence}
\begingroup\scriptsize
\begin{verbatim}
.venv/bin/python scripts/calibrate_giu_preset.py --check
\end{verbatim}
\endgroup
This verifies the local Spring 2023 source hashes, extraction aggregates, importer projection, preset payload hash, and the checked-in limitation-aware calibration artifact (SHA-256 \path{c2620b38ed4535fbbeb55ca66ec3c031cf346a1a5dc6c66d62cc545c5df53b73}).
The command verifies historical provenance only; the current replay contains 148 validation errors and does not provide schedule/model agreement or current institutional approval.
The questionnaire and acceptance gates are in \path{docs/GIU_INSTITUTIONAL_VALIDATION_PROTOCOL.md}.

\subsection*{Disruption Experiments}
\begingroup\scriptsize
\begin{verbatim}
.venv/bin/python scripts/run_disruption_experiments.py \
  --mode small_demo --seeds 1,2,3 \
  --room-mode decomposed \
  --out paper/disruption_results.jsonl
\end{verbatim}
\endgroup

\subsection*{External Formats}
\begingroup\scriptsize
\begin{verbatim}
.venv/bin/python scripts/import_external_benchmark.py \
  itc2007 instance.ctt --out instance.json
.venv/bin/python scripts/import_external_benchmark.py \
  itc2019 instance.xml --out inspection.json
\end{verbatim}
\endgroup
The second command is explicitly an inspection report, not a Planora instance conversion.
The ITC-2019 solver path is conditional student sectioning after class placements are fixed; it is not a full joint placement-and-sectioning conversion, and the local validators are not the official distribution/objective validator.

\subsection*{Checked-in ITC-2007 Evidence}
The official-validator breadth report is \path{reports/itc2007_breadth_21_rescue_seed17_2026-08-11.json}, SHA-256 \path{50adfe6ccc8add2e7e9225f8268c430331153bdb8244826eec125ad5716ea385}.
The final corrected room-dive report is \path{reports/itc2007_fixed_time_room_dive_breadth_final_v2_seed17_2026-08-11.json}, SHA-256 \path{24fccbf81d78a1371d625f42c94e1ff1b7d8e841633874ddb788a552d7146be6}.
Its counterbalanced matrix index is \path{output/itc2007-room-dive-breadth-seed17-counterbalanced-final-v2/matrix_index.json}, SHA-256 \path{aba9413209699d2eff0db431f90e971dea3dada0e675420f3d9cefeb6bc8f2f4}, with source SHA-256 \path{a63732de42d9e0053ea196f46a410024b85c006e4cc6d7f91e5529f3bd1294cb}.
That source snapshot predates the later repeat-week greedy-to-decomposed correctness-only change; retain the report and matrix as immutable historical evidence and do not treat their source-stability gate as exact-current-checkout equivalence.
The excluded failed breadth report remains at \path{reports/itc2007_fixed_time_room_dive_breadth_seed17_2026-08-11.json}, SHA-256 \path{43560980d667a5495dda0b597ae4942f4e355ed94ec041dbdb0fcd3cd2d912d7}.
The targeted seven-instance deadline-correction gate is \path{reports/itc2007_fixed_time_room_dive_deadline_gate_v2_2026-08-11.json}, SHA-256 \path{d595fb73473c804e66eceba47301cc12fb4c7522ee17e1afb75c578bd624fcc2}; the later 42-record matrix is the full corrected result.
These paths are evidence inputs, not commands that silently refetch official instances or validators.

The bounded support-proxy mechanism record is \path{paper/evidence/itc2007_support_proxy_comp10_2026-08-13.json}.
It binds the instance, retained incumbent, comparator, implementation files, validator, two byte-identical outputs, fixed work quotas, and separate post-incumbent timings.
The selected quality ledger is \path{paper/evidence/selected_post_incumbent_quality_2026-08-13.json}.
It records instance, incumbent, output, comparator, proof-report, and implementation hashes; it also marks whether each proof implementation hash still matches the checkout.
The ledger is intentionally not classified as a replicated or equal-budget benchmark result.
The record explicitly marks the experiment as mechanism evidence rather than an end-to-end solve or equal-runtime comparison.

The current end-to-end paired evidence is \path{paper/evidence/itc2007_fresh_redundant_matrix_2026-08-13.json}, SHA-256 \path{ef02acc47d02c59462a00a4ba57b25b1e39ce02a6d98dd7061c1e28cc269409e}.
It binds two complete 42-row passes to the identical Planora source, CPSolver classes, validator, instances, seed, worker count, CPU, and configured search budget, and reports pass-to-pass winner stability explicitly.
The fresh competition matrix is \path{output/itc2019-open-source-competition-30-seed17-10s-20260813/report.json}; it contains 120/120 planned rows, exact input/tool/source provenance, the non-equal-resource claim boundary, and authenticated official agreement for the three emitted complete valid outputs. The authenticated Results-page snapshot is \path{paper/evidence/itc2019_official_results_page_2026-08-13.json}. The earlier failed-login record at \path{paper/evidence/itc2019_official_validator_attempt_2026-08-13.json} is retained only as historical provenance and is superseded for access status.

\subsection*{Release Artifact Freeze}
\begingroup\small
\begin{verbatim}
./scripts/freeze_release_artifacts.sh \
  v1.0-rc-20260811-a
./scripts/verify_release_artifact.sh \
  release/v1.0-rc-20260811-a
\end{verbatim}
\endgroup
The identifier must be new: the freezer is immutable and refuses to overwrite an existing release directory. The checked-in \path{release/v1.0} predates the current evidence and must not be presented as the current release candidate.
The workflow is designed to preserve paper paths, source required by the experiment scripts, Git/environment evidence, and non-self-referential checksums.
The freezer runs the verifier before reporting success.
No newly frozen and verified bundle is claimed by this paper revision.

\subsection*{Paper Build}
\begingroup\small
\begin{verbatim}
cd paper
latexmk -pdf -interaction=nonstopmode \
  -halt-on-error main.tex
\end{verbatim}
\endgroup

Source licensing is \texttt{AGPL-3.0-or-later}.


\end{document}
